\documentclass{beamer}
%
% Choose how your presentation looks.
%
% For more themes, color themes and font themes, see:
% http://deic.uab.es/~iblanes/beamer_gallery/index_by_theme.html
%
\mode<presentation>
{
  \usetheme{Madrid}      % or try Darmstadt, Madrid, Warsaw, ...
  \usecolortheme{beaver} % or try albatross, beaver, crane, ...
  \usefonttheme{serif}  % or try serif, structurebold, ...
  \setbeamertemplate{navigation symbols}{}
  \setbeamertemplate{caption}[numbered]
} 

\usepackage[english]{babel}

\usepackage[utf8x]{inputenc}
\usepackage{xcolor}
\usepackage{listings}
\usepackage{textpos}
\usepackage{pdfpages}


\usepackage{verbatim}
\usepackage{filecontents}
\usepackage{siunitx}
\sisetup{per=slash, load=abbr}
\usepackage{tikz}
\usepackage{pgfplots}
\pgfplotsset{width=7cm,compat=newest}
\usepackage[autostyle=true,german=quotes]{csquotes}

\input{../template/custom-colors.tex}


\lstset
{
    language=[LaTeX]TeX,
    breaklines=true,
    basicstyle=\tt\scriptsize,
    %commentstyle=\color{green}
    keywordstyle=\color{blue},
    %stringstyle=\color{black}
    identifierstyle=\color{magenta},
}

\renewcommand{\rmdefault}{phv} % Arial
\renewcommand{\sfdefault}{phv} % Arial

\setbeamertemplate{frametitle}[default][center]

\addtobeamertemplate{frametitle}{}{%
\begin{textblock*}{100mm}(.88\textwidth,-.87cm)
\includegraphics[height=0.8cm,width=1.3cm]{ffhslogo}
\end{textblock*}}


\addtobeamertemplate{frametitle}{}{%
\begin{textblock*}{100mm}(-.02\textwidth,-.9cm)
\includegraphics[height=1.5cm,width=1.5cm]{chip}
\end{textblock*}}


\title[Energie in der Informatik]{Stand der Thesis}
\author{Samuel Riolo}
\institute{FFHS}
\date{18. Juni 2016}


\begin{document}

\begin{frame}
  \titlepage
\end{frame}

% Uncomment these lines for an automatically generated outline.
\begin{frame}{Outline}
\frametitle{Inhaltsverzeichnis}
    \tableofcontents[]
\end{frame}



\section{Um was geht es -- in K{\"u}rze} 
\begin{frame}
\frametitle{Um was geht es -- in K{\"u}rze} 
\begin{itemize}
\item Energieverbrauch pro Befehlssatz
\item Erstellung eines Benchmarks
\item Befehlssatz x-mal ausführen und Strom messen
\end{itemize}


\begin{figure}
  \includegraphics[width=200px,height=120px]{iot_galileo.png}
\end{figure}	

\end{frame}


\section{ToDo} 
\begin{frame}
\frametitle{ToDo}
\begin{itemize}
\item Idee und Zielsetzung (Stichworte sind vorhanden, muss jedoch noch in Sätze gefasst werden)
\item Unterkapitel: Ziele für die Energieeinsparung 2050
\item Unterkapitel: Strom-Ladung und -Entladung innerhalb eines Transistors
\item Kapitel Resultate fertigstellen
\item Fazit mit Diskussion erstellen
\end{itemize}

\end{frame}

\begin{frame}
\frametitle{Resultate in Roh­form} 
\Huge Präsentation: Diagramme im Anhang
\end{frame}

\section{Resultate}
\begin{frame}
\frametitle{Resultate} 
\begin{minipage}{0.8\textwidth}

\begin{filecontents}{dataset}
Pos	Name	Bench1	Bench2	Bench3	Avg	time	energymWh
1	nop	413	413	413	413	129147	94.2299561
2	addl\_0	410	410	409	410	145291	105.2391143333
3	addl\_32	408	408	408	408	145291	104.7257528
4	and\_0	411	412	411	411	145291	105.4957951
5	and\_32	408	408	408	408	145291	104.7257528
6	xor\_0	411	411	411	411	145291	105.4957951
7	nop\_nop	411	411	411	411	145292	105.4965212
8	subl\_0	408	408	409	408	145292	104.7264736
9	inc\_8	411	412	411	411	145292	105.4965212
10	xor\_32	410	410	410	410	161355	116.874805
11	or\_32	411	410	410	410	161434	116.9320273333
12	nop\_nop\_nop	410	410	410	410	161435	116.9327516667
13	subl\_32	411	411	411	411	161435	117.2179535
14	dec\_8	410	411	410	410	161435	116.9327516667
15	or\_0	412	412	412	412	161436	117.5038832
16	shr\_0	407	407	406	407	177578	127.6845012667
17	shr\_32	411	410	410	410	177578	128.6256646667
18	shl\_0	410	410	410	410	177578	128.6256646667
19	shl\_32	409	409	409	409	209866	151.6421760667
20	imull\_0	407	407	407	407	226008	162.5072856
21	imull\_32	407	407	407	407	226008	162.5072856
\end{filecontents}



\begin{figure}[htb] 
    \centering
    \begin{tikzpicture} 
        \begin{axis}[
            ybar,
            scale only axis,
            xlabel= Benchmarks,
            width = 0.85\textwidth,
            height = 0.5\textwidth,
            y axis line style=blau_2b!100!black,
            ylabel= Stromverbrauch in mA,
            % height = 80mm,
            ymin = 405,
            ymax = 415,
            xmin = 0,
            xmax = 22,
            bar width = 2mm,
            axis x line* = bottom,
            axis y line* = left,
            xticklabel style={rotate=90},
            xticklabels from table = {dataset}{Name},
            xtick = {1,...,21},
            bar shift = -1mm,
            y tick label style = {/pgf/number format/use comma},
            legend style = {at={(0.5, 1.025)}, anchor = south east, legend columns = -1, draw=none, area legend},
            area legend,
            extra y ticks ={405},
            extra y tick style = {grid=none, yticklabel style={yshift=2.5mm, xshift=1.9mm, rotate=0, inner sep=.5pt}},
            extra y tick labels = {\includegraphics{almost_equals.eps}}
            ]
            \addplot+[mark=none, fill=blau_2a, draw = blau_2b] table[x=Pos, y=Avg, fill=red] {dataset};
            \legend{Stromverbrauch}
        \end{axis} 
        \begin{axis}[
            ybar,
            scale only axis,
            width = 0.85\textwidth,
            height = 0.5\textwidth,
            y axis line style=gruen_4b!100!black,
            axis y line* = right,
            axis x line = none,
            ylabel = Durchlaufszeit in mS,
            % height = 80mm,
            xmin = 0,
            xmax = 22,
            ymin = 0,
            xticklabel style={rotate=90},
            xticklabels from table = {dataset}{Name},
            xtick = {1,...,21},
            ymax = 246008,
            ymin = 109147,
            bar width = 2mm,
            bar shift = 1mm,
            legend style = {at={(0.5, 1.025)}, anchor = south west, legend columns = -1, draw=none, area legend},
            area legend,
            extra y ticks ={109147},
            extra y tick style = {grid=none, yticklabel style={yshift=2.5mm, xshift=-1.6mm, rotate=0, inner sep=.5pt}},
            extra y tick labels = {\includegraphics{almost_equals.eps}},
            ]
            \addplot+[mark=none,  fill=gruen_4a, draw = gruen_4b,] table[x=Pos, y=time] {dataset};
            \legend{Durchlaufszeit}
        \end{axis}
    \end{tikzpicture}
    \caption{Ausführung aller Benchmarks auf dem Galileo Board}
    \label{fig:benchmarks_galileo}
\end{figure}


\end{minipage}
\end{frame}


\IfFileExists{../../../ResearchPoster.pdf}{%
\section{Research Poster}

\setbeamertemplate{navigation symbols}{}
\setbeamertemplate{background canvas}{\includegraphics[height = \paperheight, 
                                      width = \paperwidth, scale=0.06]{../../../ResearchPoster.pdf}}
\begin{frame}[plain]
\end{frame}
}


\end{document}

